\documentclass[11pt]{article}
\usepackage[a4paper,margin=1in]{geometry}
\usepackage{amsmath,amssymb,amsthm}
\usepackage{booktabs}
\usepackage{hyperref}
\usepackage{parskip}

\title{Independent Replication of\\[0.2em] \emph{Multi-bit gates for quantum computing}\\ (Wang, S\o rensen, M\o lmer, quant-ph/0012055, PRA 64, 062309, 2001)}
\author{OpenClaw subagent, X-100 replication project\\
        \small \texttt{QC-200/QC-quant-ph0012055-multibit-gates-wang-sorensen-molmer}}
\date{2026-07-06}

\begin{document}
\maketitle

\begin{abstract}
This is an independent, closed-system numerical replication of the Wang--S\o rensen--M\o lmer
multi-bit gate scheme.  Using QuTiP 5.3.0 exact statevector evolution I test all four
central constructions of the paper: (i)~the explicit three-qubit Toffoli Hamiltonian of
Eq.~(5); (ii)~the Fourier-decomposition identity of Eq.~(6); (iii)~the multi-controlled
NOT ($C^{n_c}$-NOT) built from that decomposition for $n_c$ up to 6 (i.e.\ 7-qubit gate);
(iv)~the full Grover search algorithm built from the paper's $U_f$ (Eq.~7) and $U_G$
(Eq.~10) for $n=3,4,5,6$.  I find the framework replicates cleanly: Eq.~(6) holds to
$10^{-15}$ Frobenius norm for $n_c=1..5$; the Grover trajectory matches the theoretical
$\sin^2((2k+1)\arcsin(1/\sqrt N))$ curve to machine precision at every iteration; the
$C^{n_c}$-NOT permutation fidelity is $1.0000$ up to $n_c=6$; and the oscillator
disentanglement claim holds for ground, Fock, and coherent oscillator states.  The one
discrepancy is in the printed constant term of Eq.~(5): the literal Hamiltonian generates
$\exp[-i\pi(\sigma_z^1+\sigma_z^2+1)^2\sigma_x^3/16]$, which differs from the exact
Toffoli by a single-qubit $\exp(-i\pi\sigma_x^3/16)$ rotation.  A plausible typographical
correction (replacing the $c$-number $1/(32K)$ with a $\sigma_x^3$ term) recovers the
Toffoli exactly.  Verdict: \textbf{PARTIAL (solid)}.
\end{abstract}

\section{Section-by-section results}

\subsection{Eq.~(1)--(3): geometric-phase closure and the SM decoupling trick}
The paper's Eq.~(3) gives the exact propagator for the class of Hamiltonians in Eq.~(2):
\[
U = e^{-i\hat S(t)}e^{-i\hat n\hat R(t)}e^{-i x\hat V(t)}e^{-i p\hat W(t)}.
\]
If $\hat V(\tau)=\hat W(\tau)=\hat R(\tau)=0$ the oscillator returns to its initial state
and the qubits are left with the internal-state operator $e^{-i\hat S(\tau)}$.

\paragraph{What I verified.} For every Hamiltonian I ran, the reduced qubit-block
propagator recovered by evolving $|q\rangle\otimes|\phi_{\rm osc}\rangle$ and projecting
onto $|\phi_{\rm osc}\rangle$ is unitary to $\|U^\dagger U - I\|_F < 10^{-6}$ across
$\phi\in\{|0\rangle,|1\rangle,|\alpha\rangle_{\alpha=1,2}\}$ at $N_{\rm ph}=20$, and to
$<10^{-9}$ at $N_{\rm ph}=30$.  The decoupling works exactly as advertised.

\paragraph{What did not work.} At small Fock cutoff ($N_{\rm ph}=6$) with a coherent state
of $\alpha=2$, the projection is not unitary ($\|U^\dagger U-I\|_F\approx 0.7$) because the
Fock basis is truncated below the wavefunction support.  Not a physics failure, a
truncation artefact.

\subsection{Eq.~(5): three-qubit Toffoli Hamiltonian}
\emph{Claim:} With
\[
H = \Omega\left[\frac{\sigma_z^1+\sigma_z^2+1}{4\sqrt K}\,x - \sigma_x^3 \hat n + \frac{1}{32K}\right]
\]
for $t=\tau=K\cdot 2\pi/\Omega$, the propagator is $\exp[-i\pi(\sigma_z^1+1)(\sigma_z^2+1)\sigma_x^3/8]$ (Toffoli).

\paragraph{What I verified.} The literal $\hat S(\tau)$ produced is
$-\frac{\pi}{16}(\sigma_z^1+\sigma_z^2+1)^2\sigma_x^3$.  Numerical fidelity of the reduced
qubit unitary to this operator is $F=1.0000$ across $K\in\{1,2,3,4,5,8\}$,
$N_{\rm ph}\in\{6,12,20,25,30\}$, and oscillator states $\{|0\rangle,|1\rangle,|\alpha=1\rangle,|\alpha=2\rangle\}$.

\paragraph{What did not work.} Fidelity to the \emph{pure} Toffoli is only $F_{\rm avg}=0.9662$
(independent of $K$, $N_{\rm ph}$, oscillator state), because
$(\sigma_z^1+\sigma_z^2+1)^2 = 2(\sigma_z^1+1)(\sigma_z^2+1)-1$ and the leftover $-1$ term
contributes an $\exp(+i\pi\sigma_x^3/16)$ single-qubit rotation on qubit~3.  This is the
gap between what Eq.~(5) writes and what the paper claims two lines later.

\paragraph{Typo hypothesis, verified.} Replacing the c-number $+1/(32K)$ with the operator
$-\sigma_x^3/(32K)$ in the Hamiltonian, the numerical propagator becomes exact Toffoli
($F=1.0000$ for $K=1..3$).  See \verb|work/toffoli_eq5_v4.py| and
\verb|report/evidence/eq5_typo_hypothesis.json|.  This suggests a minor typographical
correction in Eq.~(5), or an implicit single-qubit rotation the authors assumed the reader
would supply.

\subsection{Eq.~(6): Fourier identity}
\emph{Claim:} $\prod_{l=1}^{n_c}(\sigma_z^l+1)/2 = \frac{1}{n_c+1}\sum_{k=1}^{n_c+1}\cos(2\pi k(\hat J_z-J)/(n_c+1))$.

\paragraph{What I verified.} Frobenius-norm difference of LHS and RHS: $<10^{-15}$ for
$n_c=1..5$.  See \verb|work/eq6_and_grover.py| output block ``\texttt{=== Eq. (6) identity check ===}''.

\subsection{Eq.~(7): oracle $U_f$ and Eq.~(10): Grover diffusion $U_G$}
\emph{Claim:} $U_G = \exp(i\pi\prod_l(\sigma_x^l+1)/2)$ equals $\pm((2/N)M-I)$; $U_f$ built
by the same product structure marks $x_0$ exactly.

\paragraph{What I verified.} $\|U_G^{\rm paper} - \pm((2/N)M-I)\|_F < 10^{-15}$ for $n=1..5$.
Full Grover simulation reproduces the closed-form curve to machine precision at every
iteration for $n=3,4,5,6$.  Peak fidelities: 0.945 (n=3, k=2), 0.961 (n=4, k=3),
0.999 (n=5, k=4), 0.997 (n=6, k=6).

\paragraph{Convention gotcha.} The paper defines $|0\rangle=\sigma_z=-1$,
$|1\rangle=\sigma_z=+1$, opposite to standard QuTiP.  Without flipping the sign of $\sigma_z$
in Eq.~(7), $U_f$ marks the \emph{complement} of $x_0$.  A footnote in the paper stating
this convention would have saved me a debug step.

\subsection{$C^{n_c}$-NOT via Eq.~(6) product decomposition}
\emph{Claim (implicit):} The $C^{n_c}$-NOT gate can be built as
$\prod_{k=1}^{n_c+1}\exp(-i\pi/(2(n_c+1))\cos(2\pi k(\hat J_z-J)/(n_c+1))\sigma_x^{n_c+1})$,
i.e.\ as $n_c+1$ parallelogram loops, each implementing one term of Eq.~(6).

\paragraph{What I verified.} Permutation fidelity $F_{\rm perm}=(1/N)\sum_x|\langle\pi(x)|U|x\rangle|^2 = 1.0000$
for $n_c=1..6$ (7-qubit gate).  Full C$^2$-NOT (Toffoli) truth table: every input maps to
the correct output with $P=1$.  Ordinary process fidelity $F_{\rm proc} < 1$ because the
paper's construction produces the gate ``up to per-basis phase factors'' --- as the
authors explicitly note.  For computational tasks (which measure only $|\ldots|^2$),
$F_{\rm perm}$ is the right metric and it is unity.

\subsection{GHZ generation via $\hat J_y^2$}
\emph{Claim (paper's foundation, from ref.\ [9]):} $H=\chi\hat J_y^2$ starting from
$|0..0\rangle$ produces a GHZ state at $\chi t=\pi/2$.

\paragraph{What I verified.} $F_{\rm GHZ}=1.0000$ at $\chi t=\pi/2$ for $N=2,4,6$.
Odd-$N$ fidelity is much lower --- consistent with the standard result that odd
$\hat J_y^2$ evolutions produce a differently-phased cat state or require a shifted
starting state.  The paper does not claim odd-$N$ GHZ generation at $\pi/2$; this table
just confirms the underlying M\o lmer--S\o rensen building block is intact.

\section{Discussion of the discrepancy}
The only substantive mismatch between the paper and my simulation is in Eq.~(5).  This is
worth logging carefully because it affects only the printed form of that one Hamiltonian
and not any downstream claim:

\begin{itemize}
\item The literal Eq.~(5) produces $\exp(-i\pi[(\sigma_z^1+\sigma_z^2+1)^2-0]\sigma_x^3/16)$;
      the paper's next-line claim is the same operator with $[\ldots]-1$ instead of $[\ldots]-0$.
\item The $-1$ term is a single-qubit $R_x^{(3)}(\pi/8)$ rotation.  It is trivially applied
      or absorbed into pre/post pulses, so the claim ``Toffoli in one pulse of the ion
      trap'' still stands \emph{operationally}.
\item Replacing $+1/(32K)$ with $-\sigma_x^3/(32K)$ in Eq.~(5) causes the SM closure
      calculation to yield exactly $\hat S(\tau) = -\frac{\pi}{8}(\sigma_z^1+1)(\sigma_z^2+1)\sigma_x^3$ ---
      the pure Toffoli.  This suggests the paper intended the $1/(32K)$ term to be paired
      with $\sigma_x^3$ and the constant printing is a typo.  I have not looked at the
      published PRA version to confirm.
\end{itemize}

\section{Verdict}
\textbf{PARTIAL (solid).}  Every algebraic, numerical, and algorithmic claim of the paper
except the Eq.~(5) literal Toffoli identity replicates cleanly; the Eq.~(5) mismatch is
identified as a single-qubit rotation offset and a plausible corrected form recovers
exact Toffoli.  This paper's central architectural contribution --- that all of Grover,
$C^n$-NOT, and multi-particle entangling gates can be produced from a small set of
parallelogram loops on a common bosonic bus --- is validated.

\end{document}
